\documentclass[10pt,a4paper]{amsart}
\usepackage{amsmath,amssymb,amsthm,graphicx}
\usepackage{verbatim}
\newenvironment{reference}{\comment}{\endcomment}
%\newenvironment{reference}{}{}
\newenvironment{slogan}{\comment}{\endcomment}
\newenvironment{history}{\comment}{\endcomment}

% For commutative diagrams we use Xy-pic
\usepackage[all]{xy}

% We use 2cell for 2-commutative diagrams.
\xyoption{2cell}
% Theorem environments.
%
\theoremstyle{plain}
\newtheorem{theorem}[subsection]{Theorem}
\newtheorem{proposition}[subsection]{Proposition}
\newtheorem{lemma}[subsection]{Lemma}

\theoremstyle{definition}
\newtheorem{definition}[subsection]{Definition}
\newtheorem{example}[subsection]{Example}
\newtheorem{exercise}[subsection]{Exercise}
\newtheorem{situation}[subsection]{Situation}

\theoremstyle{remark}
\newtheorem{remark}[subsection]{Remark}
\newtheorem{remarks}[subsection]{Remarks}

\numberwithin{equation}{subsection}

% Macros
%
\def\lim{\mathop{\mathrm{lim}}\nolimits}
\def\colim{\mathop{\mathrm{colim}}\nolimits}
\def\Spec{\mathop{\mathrm{Spec}}}
\def\Hom{\mathop{\mathrm{Hom}}\nolimits}
\def\Ext{\mathop{\mathrm{Ext}}\nolimits}
\def\SheafHom{\mathop{\mathcal{H}\!\mathit{om}}\nolimits}
\def\SheafExt{\mathop{\mathcal{E}\!\mathit{xt}}\nolimits}
\def\Sch{\mathit{Sch}}
\def\Mor{\mathop{\mathrm{Mor}}\nolimits}
\def\Ob{\mathop{\mathrm{Ob}}\nolimits}
\def\Sh{\mathop{\mathit{Sh}}\nolimits}
\def\NL{\mathop{N\!L}\nolimits}
\def\CH{\mathop{\mathrm{CH}}\nolimits}
\def\proetale{{pro\text{-}\acute{e}tale}}
\def\etale{{\acute{e}tale}}
\def\QCoh{\mathit{QCoh}}
\def\Ker{\mathop{\mathrm{Ker}}}
\def\Im{\mathop{\mathrm{Im}}}
\def\Coker{\mathop{\mathrm{Coker}}}
\def\Coim{\mathop{\mathrm{Coim}}}

% Boxtimes
%
\DeclareMathSymbol{\boxtimes}{\mathbin}{AMSa}{"02}

%
% Macros for moduli stacks/spaces
%
\def\QCohstack{\mathcal{QC}\!\mathit{oh}}
\def\Cohstack{\mathcal{C}\!\mathit{oh}}
\def\Spacesstack{\mathcal{S}\!\mathit{paces}}
\def\Quotfunctor{\mathrm{Quot}}
\def\Hilbfunctor{\mathrm{Hilb}}
\def\Curvesstack{\mathcal{C}\!\mathit{urves}}
\def\Polarizedstack{\mathcal{P}\!\mathit{olarized}}
\def\Complexesstack{\mathcal{C}\!\mathit{omplexes}}
% \Pic is the operator that assigns to X its picard group, usage \Pic(X)
% \Picardstack_{X/B} denotes the Picard stack of X over B
% \Picardfunctor_{X/B} denotes the Picard functor of X over B
\def\Pic{\mathop{\mathrm{Pic}}\nolimits}
\def\Picardstack{\mathcal{P}\!\mathit{ic}}
\def\Picardfunctor{\mathrm{Pic}}
\def\Deformationcategory{\mathcal{D}\!\mathit{ef}}
\usepackage{cleveref}
\crefname{originaldefinition}{original definition}{original definitions}
\Crefname{originaldefinition}{Original definition}{Original definitions}
\makeatletter
\newlabel{definition-separated}{{00WA}{source}{}{sourcepointer-1}{}}
\newlabel{definition-separated@cref}{{[originaldefinition][0][]00WA}{[1][0][]source}{}{}{}}
\makeatother
\crefname{originallemma}{original lemma}{original lemmas}
\Crefname{originallemma}{Original lemma}{Original lemmas}
\makeatletter
\newlabel{lemma-colimit-sheaves}{{00WI}{source}{}{sourcepointer-2}{}}
\newlabel{lemma-colimit-sheaves@cref}{{[originallemma][0][]00WI}{[1][0][]source}{}{}{}}
\makeatother
\crefname{originaltheorem}{original theorem}{original theorems}
\Crefname{originaltheorem}{Original theorem}{Original theorems}
\makeatletter
\newlabel{theorem-plus}{{00WB}{source}{}{sourcepointer-3}{}}
\newlabel{theorem-plus@cref}{{[originaltheorem][0][]00WB}{[1][0][]source}{}{}{}}
\makeatother
\crefname{originallemma}{original lemma}{original lemmas}
\Crefname{originallemma}{Original lemma}{Original lemmas}
\makeatletter
\newlabel{categories-lemma-directed-commutes}{{002W}{source}{}{sourcepointer-4}{}}
\newlabel{categories-lemma-directed-commutes@cref}{{[originallemma][0][]002W}{[1][0][]source}{}{}{}}
\makeatother
\crefname{originallemma}{original lemma}{original lemmas}
\Crefname{originallemma}{Original lemma}{Original lemmas}
\makeatletter
\newlabel{lemma-image-of-quasi-compact}{{0GMP}{source}{}{sourcepointer-5}{}}
\newlabel{lemma-image-of-quasi-compact@cref}{{[originallemma][0][]0GMP}{[1][0][]source}{}{}{}}
\makeatother
\title{Global sections of a filtered colimit of sheaves of sets}
\author{The Stacks Project Authors}
\date{stacks-project:sites.tex:local-snapshot-20261012; GNU FDL 1.2 or later}
\begin{document}
\maketitle
\noindent\textit{Editorial note.} All delivered mathematical text is the author's own, copied from the pinned Stacks Project chapter file sites.tex; no statement or proof step was rewritten, reordered or removed.\par\medskip
\noindent\textit{Editorial note.} The proof is the author's complete argument for all four parts, including the equalizer argument in part (4). The author's sentence 'we urge the reader to skip the rest of the proof' (source line 3844-3845) introduces a reading suggestion and the argument is nevertheless written out in full immediately afterwards; it is preserved verbatim inside the counted span.\par\medskip
\noindent\textit{Editorial note.} Every labelled reference occurring inside the counted spans is carried in assembly.external\_references as a non-load-bearing navigation pointer (same-chapter ones at their definition site, cross-chapter ones at the invocation inside the counted span), so no load-bearing mathematics is replaced by a stub. No bibliography citation occurs inside the counted spans.\par\medskip
\noindent\textit{Editorial note.} The checked chapter list for same-name collisions is the set of stacks chapters staged in this repository: chow, cohomology, divisors, etale-cohomology, flat, more-algebra, more-morphisms, morphisms, spaces-more-morphisms, varieties together with this card's chapter sites. The label occurs in none of them, so no cross-chapter restatement had to be weighed.\par\medskip
\noindent\textit{Editorial note.} The preamble is the chapter's own native preamble (sources/stacks-sites/preamble.tex): the comment-type definitions of reference, slogan and history (lines 8-13), the Xy-pic setup (lines 14-18) and the native theorem declarations with the mathematical macros (lines 154 to the file's actual last line).\par\medskip
\noindent\textit{Allowed standard prerequisites (editorial).}\begin{itemize}
\item Separated presheaves and the plus construction. Same chapter, Definition definition-separated and Theorem theorem-plus, together with Lemma lemma-colimit-sheaves describing the sheafification of a filtered colimit. Supplies the injectivity of the map from a separated presheaf to its sheafification used in part (1).
\item Filtered colimits of subobjects of the final object. Categories, Lemma categories-lemma-directed-commutes. Identifies the equalizer of the colimit maps with the colimit of the equalizers.
\item Quasi-compactness of the final sheaf and of the topos. Same chapter, the notion of a quasi-compact object and Lemma lemma-image-of-quasi-compact. Turns the surjectivity of the displayed maps onto the final object into the existence of a stage where the equalizer is everything.
\end{itemize}\medskip
\ifdefined\hypertarget\hypertarget{sourcepointer-1}{}\fi
\noindent\textit{Source reference.} Original definition 00WA: sources/stacks-sites/pinned-original-source.tex, source lines 3766--3766. Navigation pointer invoked at proof line 3766; the counted span contains the pointer only and the cited result is not reproduced inside it, so the pointer is registered as a non-load-bearing navigation stub in the established form for Stacks material. No bibliography citation is involved.\par\medskip
\ifdefined\hypertarget\hypertarget{sourcepointer-2}{}\fi
\noindent\textit{Source reference.} Original lemma 00WI: sources/stacks-sites/pinned-original-source.tex, source lines 3767--3767. Navigation pointer invoked at proof line 3767; the counted span contains the pointer only and the cited result is not reproduced inside it, so the pointer is registered as a non-load-bearing navigation stub in the established form for Stacks material. No bibliography citation is involved.\par\medskip
\ifdefined\hypertarget\hypertarget{sourcepointer-3}{}\fi
\noindent\textit{Source reference.} Original theorem 00WB: sources/stacks-sites/pinned-original-source.tex, source lines 3673--3673. Navigation pointer invoked at proof line 3673; the counted span contains the pointer only and the cited result is not reproduced inside it, so the pointer is registered as a non-load-bearing navigation stub in the established form for Stacks material. No bibliography citation is involved.\par\medskip
\ifdefined\hypertarget\hypertarget{sourcepointer-4}{}\fi
\noindent\textit{Source reference.} Original lemma 002W: sources/stacks-sites/pinned-original-source.tex, source lines 3788--3788. Navigation pointer invoked at proof line 3788; the counted span contains the pointer only and the cited result is not reproduced inside it, so the pointer is registered as a non-load-bearing navigation stub in the established form for Stacks material. No bibliography citation is involved.\par\medskip
\ifdefined\hypertarget\hypertarget{sourcepointer-5}{}\fi
\noindent\textit{Source reference.} Original lemma 0GMP: sources/stacks-sites/pinned-original-source.tex, source lines 3810--3810. Navigation pointer invoked at proof line 3810; the counted span contains the pointer only and the cited result is not reproduced inside it, so the pointer is registered as a non-load-bearing navigation stub in the established form for Stacks material. No bibliography citation is involved.\par\medskip
\begin{lemma}
\label{lemma-directed-colimits-global-sections}
Let $\mathcal{C}$ be a site. Let
$\mathcal{I} \to \Sh(\mathcal{C})$, $i \mapsto \mathcal{F}_i$
be a filtered diagram of sheaves of sets.
Consider the canonical map
$$
\Psi :
\colim_i \Gamma(\mathcal{C}, \mathcal{F}_i)
\longrightarrow
\Gamma(\mathcal{C}, \colim_i \mathcal{F}_i)
$$
We have the following:
\begin{enumerate}
\item If all the transition maps are injective then $\Psi$ is injective.
\item If $\Sh(\mathcal{C})$ is quasi-compact, then $\Psi$ is injective.
\item If $\Sh(\mathcal{C})$ is quasi-compact and all the transition maps
are injective then $\Psi$ is an isomorphism.
\item Assume there exists a set $S \subset \Ob(\Sh(\mathcal{C}))$
with the following properties:
\begin{enumerate}
\item for every surjection $\mathcal{F} \to *$ there exists a
$\mathcal{K} \in S$ and a map $\mathcal{K} \to \mathcal{F}$
such that $\mathcal{K} \to *$ is surjective,
\item for $\mathcal{K} \in S$ the product
$\mathcal{K} \times \mathcal{K}$ is quasi-compact.
\end{enumerate}
Then $\Psi$ is bijective.
\end{enumerate}
\end{lemma}

\begin{proof}
Proof of (1). Assume all the transition maps are injective.
In this case the presheaf
$\mathcal{F}' : V \mapsto \colim_i \mathcal{F}_i(V)$ is
separated (see Definition \ref{definition-separated}).
By Lemma \ref{lemma-colimit-sheaves}
we have
$(\mathcal{F}')^\# = \colim_i \mathcal{F}_i$.
By Theorem \ref{theorem-plus}
we see that $\mathcal{F}' \to (\mathcal{F}')^\#$ is injective.
This proves (1).

\medskip\noindent
Proof of (2). Assume $\Sh(\mathcal{C})$ is quasi-compact. Recall that
$\Gamma(\mathcal{C}, \mathcal{F}) = \Mor(*, \mathcal{F})$
for all $\mathcal{F}$ in $\Sh(\mathcal{C})$.
Let $a_i, b_i : * \to \mathcal{F}_i$ and for
$i' \geq i$ denote $a_{i'}, b_{i'} : * \to \mathcal{F}_{i'}$
the composition with the transition maps of the system.
Set $a = \colim_{i' \geq i} a_{i'}$ and similarly for $b$.
For $i' \geq i$ denote
$$
E_{i'} = \text{Equalizer}(a_{i'}, b_{i'}) \subset *
\quad\text{and}\quad
E = \text{Equalizer}(a, b) \subset *
$$
By Categories, Lemma \ref{categories-lemma-directed-commutes} we have
$E = \colim_{i' \geq i} E_{i'}$. It follows that
$\coprod_{i' \geq i} E_{i'} \to E$
is a surjective map of sheaves. Hence, if $E = *$, i.e., if $a = b$,
then because $*$ is quasi-compact, we see that $E_{i'} = *$ for some
$i' \geq i$, and we conclude $a_{i'} = b_{i'}$ for some $i' \geq i$.
This proves (2).

\medskip\noindent
Proof of (3). Assume $\Sh(\mathcal{C})$ is quasi-compact and all
transition maps are injective. Let $a : * \to \colim \mathcal{F}_i$
be a map. Then $E_i = a^{-1}(\mathcal{F}_i) \subset *$ is a subsheaf
and we have $\colim E_i = *$ (by the reference above). Hence for
some $i$ we have $E_i = *$ and we see that the image of $a$ is
contained in $\mathcal{F}_i$ as desired.

\medskip\noindent
Proof of (4). Let $S \subset \Ob(\Sh(\mathcal{C}))$ satisfy (4)(a), (b).
Applying (4)(a) to $\text{id} : * \to *$ we find there exists a
$\mathcal{K} \in S$ such that $\mathcal{K} \to *$ is surjective.
The maps $\mathcal{K} \times \mathcal{K} \to \mathcal{K} \to *$
are surjective.
By (4)(b) and Lemma \ref{lemma-image-of-quasi-compact} we conclude
that $\mathcal{K}$ and $\Sh(\mathcal{C})$ are quasi-compact.
Thus $\Psi$ is injective by (2).
Set $\mathcal{F} = \colim \mathcal{F}_i$.
Let $s : * \to \mathcal{F}$ be a global section of the colimit.
Since $\coprod \mathcal{F}_i \to \mathcal{F}$ is surjective,
we see that the projection
$$
\coprod\nolimits_{i \in I} * \times_{s, \mathcal{F}} \mathcal{F}_i \to *
$$
is surjective. By (4)(a) we obtain $\mathcal{K} \in S$ and a map
$\mathcal{K} \to \coprod_{i \in I} * \times_{s, \mathcal{F}} \mathcal{F}_i$
with $\mathcal{K} \to *$ surjective.
Since $\mathcal{K}$ is quasi-compact we obtain a factorization
$\mathcal{K} \to \coprod_{i' \in I'} * \times_{s, \mathcal{F}} \mathcal{F}_{i'}$
for some finite subset $I' \subset I$. Let $i \in I$ be an upper
bound for the finite subset $I'$. The transition maps define a map
$\coprod_{i' \in I'} \mathcal{F}_{i'} \to \mathcal{F}_i$.
This in turn produces a map
$\mathcal{K} \to * \times_{s, \mathcal{F}} \mathcal{F}_i$.
In other words, we obtain $\mathcal{K} \in S$ with
$\mathcal{K} \to *$ surjective and a commutative diagram
$$
\xymatrix{
\mathcal{K} \times \mathcal{K} \ar@<1ex>[r] \ar@<-1ex>[r] &
\mathcal{K} \ar[d] \ar[r] & {*} \ar[d]^s \\
& \mathcal{F}_i \ar[r] & \mathcal{F} \ar@{=}[r] & \colim \mathcal{F}_i
}
$$
Observe that the top row of this diagram is a coequalizer.
Hence it suffices to show that after increasing $i$ the two
induced maps
$a_i, b_i : \mathcal{K} \times \mathcal{K} \to \mathcal{F}_i$
are equal. This is done shown in the next paragraph
using the exact same argument as in the proof of (2)
and we urge the reader to skip the rest of the proof.

\medskip\noindent
For $i' \geq i$ denote
$a_{i'}, b_{i'} : \mathcal{K} \times \mathcal{K} \to \mathcal{F}_{i'}$
the composition of $a_i, b_i$ with the transition maps of the system.
Set $a = \colim_{i' \geq i} a_{i'} :
\mathcal{K} \times \mathcal{K} \to \mathcal{F}$ and similarly for $b$.
We have $a = b$ by the commutativity of the diagram above.
For $i' \geq i$ denote
$$
E_{i'} = \text{Equalizer}(a_{i'}, b_{i'}) \subset
\mathcal{K} \times \mathcal{K}
\quad\text{and}\quad
E = \text{Equalizer}(a, b) \subset
\mathcal{K} \times \mathcal{K}
$$
By Categories, Lemma \ref{categories-lemma-directed-commutes} we have
$E = \colim_{i' \geq i} E_{i'}$. It follows that
$\coprod_{i' \geq i} E_{i'} \to E$
is a surjective map of sheaves. Since $a = b$ we have
$E = \mathcal{K} \times \mathcal{K}$.
As $\mathcal{K} \times \mathcal{K}$ is quasi-compact by (4)(b),
we see that $E_{i'} = \mathcal{K} \times \mathcal{K}$ for some
$i' \geq i$, and we conclude $a_{i'} = b_{i'}$ for some $i' \geq i$.
\end{proof}

\end{document}
